
\documentclass{article}
\usepackage{colortbl}
\usepackage{makecell}
\usepackage{multirow}
\usepackage{supertabular}

\begin{document}

\newcounter{utterance}

\centering \large Interaction Transcript for game `dond', experiment `coop', episode 6 with Qwen3.6{-}35B{-}A3B{-}FP8{-}without{-}reasoning.
\vspace{24pt}

{ \footnotesize  \setcounter{utterance}{1}
\setlength{\tabcolsep}{0pt}
\begin{supertabular}{c@{$\;$}|p{.15\linewidth}@{}p{.15\linewidth}p{.15\linewidth}p{.15\linewidth}p{.15\linewidth}p{.15\linewidth}}
    \# & $\;$A & \multicolumn{4}{c}{Game Master} & $\;\:$B\\
    \hline

    \theutterance \stepcounter{utterance}  
    & & \multicolumn{4}{p{0.6\linewidth}}{
        \cellcolor[rgb]{0.9,0.9,0.9}{
            \makecell[{{p{\linewidth}}}]{
                \texttt{\tiny{[P1$\langle$GM]}}
                \texttt{Mängite koostöö läbirääkimiste mängu, milles peate saavutama teise mängijaga kokkulepet selle kohta, kuidas objektide komplekti teie vahel jagada.} \\
\\ 
\texttt{Reeglid:} \\
\texttt{(a) Teie ja teine mängija saate objektide komplekti. Igaühele antakse ka salajane väärtuse funktsioon, mis näitab, kui palju hindasite iga objekti tüüpi.} \\
\texttt{(b) Vahetate teise mängijaga sõnumeid, et kokku leppida selles, kes saab milliseid objekte. Igaüks võib saata kuni 5 sõnumit või lõpetada mängu varajalt, tehes igal ajal salajase ettepaneku.} \\
\texttt{(c) Teil palutakse esitada salajane ettepanek, mis näitab objekte, mida soovite, vormindatuna ruuduliste sulgudega järgmiselt: "{[}Ettepanek: <arv> <objekti nimi>, <arv> <objekti nimi>, <...>{]}"} \\
\texttt{(d) Kui teie ettepanekud on täiendavad, st ettepaanekute täitmiseks on piisavalt objekte, saab iga mängija skoori vastavalt tema saanud objektide väärtuste summale. Muul juhul saavad mõlemad nullpunkte.} \\
\texttt{(e) Teie eesmärk on maksimeerida oma skoori ja teise mängija skoori summa.} \\
\\ 
\texttt{Alustame.} \\
\\ 
\texttt{Saadaolevate objektide komplekt on:} \\
\\ 
\texttt{2 tangid, 1 pudel, 3 lambid, 1 kartul} \\
\\ 
\texttt{Teie salajane väärtuse funktsioon on:} \\
\\ 
\texttt{tangid: 4, pudel: 2, lamp: 0, kartul: 0} \\
\\ 
\texttt{OLULINE: Teie sõnumid, välja arvatud salajane ettepanek, edastatakse otse teisele mängijale, nii et ärge lisage mingit vastust reeglitele ega teksti, mis teie sõnumi kuulutaks. Salajase ettepaneku tegemiseks kasutage määratletud vormingut. Ärge kasutage ruudulisi sulge, kui suhtlesite teise mängijaga, või seda interpreteeritakse teie salajase ettepanekuna.} \\
            }
        }
    }
    & \\ \\

    \theutterance \stepcounter{utterance}  
    & \multicolumn{4}{p{0.6\linewidth}}{
        \cellcolor[rgb]{0.9,0.9,0.9}{
            \makecell[{{p{\linewidth}}}]{
                \texttt{\tiny{[P1$\rangle$GM]}}
                \texttt{Tere! Mul on hea meel koostööd teha.} \\
\\ 
\texttt{Minu hinnangute põhjal on mulle kõige väärtuslikumad tangid (4 punkti) ja pudel (2 punkti). Lambid ja kartul on mulle väärtusetud (0 punkti).} \\
\\ 
\texttt{Kuna meie eesmärk on maksimeerida mõlema mängija skooride summat, soovitaksin jagada objekte järgmiselt:} \\
\texttt{{-} Mina saan **2 tangit** ja **1 pudeli**.} \\
\texttt{{-} Te saate **3 lambit** ja **1 kartuli**.} \\
\\ 
\texttt{See jaotus annab mulle 6 punkti. Kui teile on lambid ja kartul väärtuslikud, saate te ka oma skoori maksimeerida. Kui teile on need väärtusetud, saate te siiski midagi, mis võib olla teile kasulik või mida saate edasi vahetada.} \\
\\ 
\texttt{Kas olete selle jaotusega nõus? Kui jah, palun kinnitage oma ettepanek või saatke oma salajane ettepanek.} \\
            }
        }
    }
    & & \\ \\

    \theutterance \stepcounter{utterance}  
    & & & \multicolumn{2}{p{0.3\linewidth}}{
        \cellcolor[rgb]{0.9,0.9,0.9}{
            \makecell[{{p{\linewidth}}}]{
                \texttt{\tiny{[GM$|$GM]}}
                \texttt{continue} \\
            }
        }
    }
    & & \\ \\

    \theutterance \stepcounter{utterance}  
    & & & \multicolumn{2}{p{0.3\linewidth}}{
        \cellcolor[rgb]{0.9,0.9,0.9}{
            \makecell[{{p{\linewidth}}}]{
                \texttt{\tiny{[GM$|$GM]}}
                \texttt{Tere! Mul on hea meel koostööd teha.} \\
\\ 
\texttt{Minu hinnangute põhjal on mulle kõige väärtuslikumad tangid (4 punkti) ja pudel (2 punkti). Lambid ja kartul on mulle väärtusetud (0 punkti).} \\
\\ 
\texttt{Kuna meie eesmärk on maksimeerida mõlema mängija skooride summat, soovitaksin jagada objekte järgmiselt:} \\
\texttt{{-} Mina saan **2 tangit** ja **1 pudeli**.} \\
\texttt{{-} Te saate **3 lambit** ja **1 kartuli**.} \\
\\ 
\texttt{See jaotus annab mulle 6 punkti. Kui teile on lambid ja kartul väärtuslikud, saate te ka oma skoori maksimeerida. Kui teile on need väärtusetud, saate te siiski midagi, mis võib olla teile kasulik või mida saate edasi vahetada.} \\
\\ 
\texttt{Kas olete selle jaotusega nõus? Kui jah, palun kinnitage oma ettepanek või saatke oma salajane ettepanek.} \\
            }
        }
    }
    & & \\ \\

    \theutterance \stepcounter{utterance}  
    & & \multicolumn{4}{p{0.6\linewidth}}{
        \cellcolor[rgb]{0.9,0.9,0.9}{
            \makecell[{{p{\linewidth}}}]{
                \texttt{\tiny{[GM$\rangle$P2]}}
                \texttt{Mängite koostöö läbirääkimiste mängu, milles peate saavutama teise mängijaga kokkulepet selle kohta, kuidas objektide komplekti teie vahel jagada.} \\
\\ 
\texttt{Reeglid:} \\
\texttt{(a) Teie ja teine mängija saate objektide komplekti. Igaühele antakse ka salajane väärtuse funktsioon, mis näitab, kui palju hindasite iga objekti tüüpi.} \\
\texttt{(b) Vahetate teise mängijaga sõnumeid, et kokku leppida selles, kes saab milliseid objekte. Igaüks võib saata kuni 5 sõnumit või lõpetada mängu varajalt, tehes igal ajal salajase ettepaneku.} \\
\texttt{(c) Teil palutakse esitada salajane ettepanek, mis näitab objekte, mida soovite, vormindatuna ruuduliste sulgudega järgmiselt: "{[}Ettepanek: <arv> <objekti nimi>, <arv> <objekti nimi>, <...>{]}"} \\
\texttt{(d) Kui teie ettepanekud on täiendavad, st ettepaanekute täitmiseks on piisavalt objekte, saab iga mängija skoori vastavalt tema saanud objektide väärtuste summale. Muul juhul saavad mõlemad nullpunkte.} \\
\texttt{(e) Teie eesmärk on maksimeerida oma skoori ja teise mängija skoori summa.} \\
\\ 
\texttt{Alustame.} \\
\\ 
\texttt{Saadaolevate objektide komplekt on:} \\
\\ 
\texttt{2 tangid, 1 pudel, 3 lambid, 1 kartul} \\
\\ 
\texttt{Teie salajane väärtuse funktsioon on:} \\
\\ 
\texttt{tangid: 2, pudel: 2, lamp: 1, kartul: 1} \\
\\ 
\texttt{OLULINE: Teie sõnumid, välja arvatud salajane ettepanek, edastatakse otse teisele mängijale, nii et ärge lisage mingit vastust reeglitele ega teksti, mis teie sõnumi kuulutaks. Salajase ettepaneku tegemiseks kasutage määratletud vormingut. Ärge kasutage ruudulisi sulge, kui suhtlesite teise mängijaga, või seda interpreteeritakse teie salajase ettepanekuna.} \\
\\ 
\texttt{Tere! Mul on hea meel koostööd teha.} \\
\\ 
\texttt{Minu hinnangute põhjal on mulle kõige väärtuslikumad tangid (4 punkti) ja pudel (2 punkti). Lambid ja kartul on mulle väärtusetud (0 punkti).} \\
\\ 
\texttt{Kuna meie eesmärk on maksimeerida mõlema mängija skooride summat, soovitaksin jagada objekte järgmiselt:} \\
\texttt{{-} Mina saan **2 tangit** ja **1 pudeli**.} \\
\texttt{{-} Te saate **3 lambit** ja **1 kartuli**.} \\
\\ 
\texttt{See jaotus annab mulle 6 punkti. Kui teile on lambid ja kartul väärtuslikud, saate te ka oma skoori maksimeerida. Kui teile on need väärtusetud, saate te siiski midagi, mis võib olla teile kasulik või mida saate edasi vahetada.} \\
\\ 
\texttt{Kas olete selle jaotusega nõus? Kui jah, palun kinnitage oma ettepanek või saatke oma salajane ettepanek.} \\
            }
        }
    }
    & \\ \\

    \theutterance \stepcounter{utterance}  
    & & & \multicolumn{4}{p{0.6\linewidth}}{
        \cellcolor[rgb]{0.9,0.9,0.9}{
            \makecell[{{p{\linewidth}}}]{
                \texttt{\tiny{[GM$\langle$P2]}}
                \texttt{Tere!} \\
\\ 
\texttt{Tundub, et meil on siin väike arusaamatus minu hinnangute osas. Minu salajane väärtuse funktsioon on järgmine:} \\
\texttt{{-} tangid: 2} \\
\texttt{{-} pudel: 2} \\
\texttt{{-} lamp: 1} \\
\texttt{{-} kartul: 1} \\
\\ 
\texttt{Seega minu hinnangul on koguväärtus: 2+2+1+1 = 6 punkti.} \\
\\ 
\texttt{Teie ettepanekus andsite mulle 2 tangit ja 1 pudeli. Minu hinnangul on see väärtus 2*2 + 2 = 6 punkti. See on minu jaoks maksimaalne võimalik skoor, kuna ma ei saa rohkem objekte.} \\
\\ 
\texttt{Teie poolt pakutud jaotus on mulle väga soodne. Kui teie hinnangute põhjal on lambid ja kartul teile väärtuslikud, siis see jaotus maksimeerib ka teie potentsiaalse skoori (eeldusel, et te hindate neid kõrgemalt kui ma).} \\
\\ 
\texttt{Kuna minu eesmärk on maksimeerida mõlema skooride summat, ja see jaotus annab mulle maksimaalse võimaliku skoori (6), olen ma selle jaotusega nõus.} \\
\\ 
\texttt{Siin on minu salajane ettepanek:} \\
\\ 
\texttt{{[}Ettepanek: 2 tangid, 1 pudel{]}} \\
            }
        }
    }
     \\ \\

    \theutterance \stepcounter{utterance}  
    & & & \multicolumn{2}{p{0.3\linewidth}}{
        \cellcolor[rgb]{0.9,0.9,0.9}{
            \makecell[{{p{\linewidth}}}]{
                \texttt{\tiny{[GM$|$GM]}}
                \texttt{proposal} \\
            }
        }
    }
    & & \\ \\

    \theutterance \stepcounter{utterance}  
    & & & \multicolumn{2}{p{0.3\linewidth}}{
        \cellcolor[rgb]{0.9,0.9,0.9}{
            \makecell[{{p{\linewidth}}}]{
                \texttt{\tiny{[GM$|$GM]}}
                [2, 1, 0, 0]
            }
        }
    }
    & & \\ \\

    \theutterance \stepcounter{utterance}  
    & & \multicolumn{4}{p{0.6\linewidth}}{
        \cellcolor[rgb]{0.9,0.9,0.9}{
            \makecell[{{p{\linewidth}}}]{
                \texttt{\tiny{[P1$\langle$GM]}}
                \texttt{Teine mängija on esitanud salajase ettepaneku. Nüüd on teie kord esitada oma ettepanek.} \\
            }
        }
    }
    & \\ \\

    \theutterance \stepcounter{utterance}  
    & \multicolumn{4}{p{0.6\linewidth}}{
        \cellcolor[rgb]{0.9,0.9,0.9}{
            \makecell[{{p{\linewidth}}}]{
                \texttt{\tiny{[P1$\rangle$GM]}}
                \texttt{{[}Ettepanek: 2 tangid, 1 pudel{]}} \\
            }
        }
    }
    & & \\ \\

    \theutterance \stepcounter{utterance}  
    & & & \multicolumn{2}{p{0.3\linewidth}}{
        \cellcolor[rgb]{0.9,0.9,0.9}{
            \makecell[{{p{\linewidth}}}]{
                \texttt{\tiny{[GM$|$GM]}}
                \texttt{proposal} \\
            }
        }
    }
    & & \\ \\

    \theutterance \stepcounter{utterance}  
    & & & \multicolumn{2}{p{0.3\linewidth}}{
        \cellcolor[rgb]{0.9,0.9,0.9}{
            \makecell[{{p{\linewidth}}}]{
                \texttt{\tiny{[GM$|$GM]}}
                [2, 1, 0, 0]
            }
        }
    }
    & & \\ \\

    \theutterance \stepcounter{utterance}  
    & & & \multicolumn{2}{p{0.3\linewidth}}{
        \cellcolor[rgb]{0.9,0.9,0.9}{
            \makecell[{{p{\linewidth}}}]{
                \texttt{\tiny{[GM$|$GM]}}
                [[2, 1, 0, 0], [2, 1, 0, 0]]
            }
        }
    }
    & & \\ \\

\end{supertabular}
}

\end{document}
